% GENERATED, edit the French unit or the translation table
\mxfront{seconde-front}
\mxcover{MATHEMATICS}{Seconde}{Annual course booklet}{15 sequences}{2026--2027}
\section*{Organisation of the booklet}
This booklet follows the fifteen stages of the corrected annual plan. The definitions, properties and methods form the course; the examples and proofs are worked in class. The gridded areas are for writing up investigations and solutions. An exercise book can supplement these spaces.
\begin{center}\small\mxarraystretch{1.25}\begin{tabularx}{\linewidth}{c X l}
\textbf{No.}&\textbf{Sequence}&\textbf{Indicative period}\\\hline
1&Arithmetic&2--11 September 2026\\
2&Vectors in the plane&14--30 September\\
3&Straight lines in the plane&1--16 October\\
4&Real numbers and intervals&2--13 November\\
5&Term review and remediation&16--20 November\\
6&Functions: representations and reference functions&23 November--4 December\\
7&Descriptive statistics of one variable&7--18 December\\
8&Algebra, equations and inequalities&4--20 January 2027\\
9&Proportions and percentage change&21--29 January\\
10&Variation and extrema of a function&1--12 February\\
11&Two-way tables of two categorical variables&1--12 March\\
12&Coordinate geometry: synthesis of vectors / lines&15--26 March\\
13&Synthesis problems and reasoning&30 March--9 April\\
14&Probability: conditioning and weighted trees&26 April--24 May\\
15&Consolidation and transition to Première&25 May--11 June\\\hline
\end{tabularx}\end{center}
The dates are progression markers, not a list of teaching days. Public holidays and school holidays are excluded from the sessions, in particular 12 March, 29 March, 6--7 May and 17 May.
\begin{keypoint}[Annual thread]
Numerical and algebraic calculation, percentages, affine functions, graph reading, reasoning and Python are revisited regularly. There is no initial revision chapter: the fluency exercises are reactivated from the first sequence.
\end{keypoint}
\begin{teachernote}The corrected DOCX plan is the reference for order and dates. Reviews 5, 12, 13 and 15 are times for reinvestment. Solutions should only be revealed after an investigation. Exercises explicitly described as fictitious are not real survey data.\end{teachernote}
\tableofcontents
